\chapter{Discussion}
\label{ch:discussion}

As mentioned in \cref{ch:methodology}, the work listed in this thesis happened in the context of easy access to vast music datasets. This access certainly influenced both the methods I used for each paper, and the research questions themselves. One can envision a quasi-mirror-image research process where instead of beginning with a research question and collecting data to find an answer, one must ask \mbox{\textit{\q{what knowledge can I pull from the available data?}}}. 

It was precisely this setting that made the combinatorial philosophical approach detailed in \cref{ch:methodology} so suitable. By immersing myself in the world of stereo image, the research questions germinated, evolved, and eventually coalesced into their current form. The need to consider what kind of knowledge was hiding in the data spurred a downstream need to understand what musical knowledge the producers and musicians could have imparted via their creative process. \rqa{} was the outcome of that thinking. By contrast, \rqb{} was a relatively straightforward concoction once \rqa{} started to materialize. One could even consider \rqb{} to be part of the answer to \rqa.


The fact that I focused on stereo image was somewhat of a practical choice. Firstly, it had a handful existing industrial use cases but potential for more. Secondly it was relatively under-studied in previous literature - certainly not at the scale that is possible today.

With that said, we must now address the proverbial elephant in the room. Despite what one might imagine from the title of this thesis, the included papers do not deal with \q{vast repositories of recorded music} directly. Rather, they all focus on understanding what is possible to find in those datasets, and subsequently how to find it. I do not consider this a contradiction in terms, but rather an optimistic indication that answers to the two minor \glspl{rq} contribute to answering the overarching \gls{rq} as well as exposing new knowledge gaps that require more research.

Previous research into stereo field generally followed a simple pattern: A novel algorithm is proposed, calculated on a mid-size music collection, and then used to predict labels that were known a-prori. The gap I identified was a lack of understanding into what musical features the existing algorithms have the ability to recognize. Once that is established, it would be possible to examine the available datasets with the intention to understand how a huge population of creatives choose to express themselves though their music. That research will reveal the scale and variation of stereo techniques that musicians are able to employ in their work. Given the basic knowledge gained from the work in this thesis, it might then be possible to reason about how musicians and producers sculpt their stereo fields, relating it both to the musical content and the wider context in which it was created.


\section{Future work}
\label{sec:future_work}

The work documented in this licentiate thesis provides the theoretical and practical basis for analysis of stereo image in music recordings. The focus was on amassing the signal processing techniques that enable the analysis of stereo field, and testing their applicability to various types of recordings that were known a-priori. The next stage of research will involve applying the knowledge from this thesis on much larger musical corpora. The signal processing techniques will be honed and run on vast musical datasets, secure in the knowledge that they will generate data appropriate to the experiment at hand.

During the licentiate, several points arose that provoked further investigation. Firstly, the available public datasets of recorded music are generally encoded using lossy compression. While lossy compression algorithms undergo perceptual tests in their development, it is unclear how they affect spatial cues at the scale of said datasets (which encompass much musical variation). The spatial effects of lossy compression therefore warrant further investigation. Secondly, the public datasets are generally formed from online services where musicians have uploaded their own recordings for free. In contrast, professional recordings (which are much more numerous) are not included in any public datasets. Since professional recordings make up the vast majority of recorded music, it needs to be assessed whether public datasets are representative of professional recordings.

In short, the doctoral research will tackle a larger problem of assessing human musical creativity though the recorded output. The techniques surfaced by the licentiate will provide the foundation, and will be applied to real-world data.